% ============================================================
% PREAMBLE - packages, page geometry, styling, helper macros
% Engine: pdfLaTeX (IPA rendered through the `tipa` package).
% Build:  pdflatex -> bibtex -> pdflatex -> pdflatex
% ============================================================

% ---------- Encoding and fonts ----------
\usepackage[T1]{fontenc}
\usepackage[utf8]{inputenc}
\usepackage{lmodern}
\usepackage{textcomp}
\usepackage{tipa}          % International Phonetic Alphabet glyphs

% ---------- Page geometry (typical binding margins) ----------
\usepackage[a4paper,left=1.5in,right=1in,top=1in,bottom=1in]{geometry}

% ---------- Spacing ----------
\usepackage{setspace}
\onehalfspacing
\setlength{\parindent}{0pt}
\setlength{\parskip}{6pt plus 2pt minus 1pt}

% ---------- Maths, tables, graphics ----------
\usepackage{amsmath,amssymb}
\usepackage{graphicx}
\usepackage{booktabs}
\usepackage{multirow}
\usepackage{array}
\usepackage{longtable}
\usepackage{tabularx}
\usepackage{caption}
\usepackage{float}
\usepackage{enumitem}
\usepackage{xcolor}

% ---------- Diagrams and plots (self-contained; no external images) ----------
\usepackage{tikz}
\usetikzlibrary{shapes.geometric,arrows.meta,positioning,fit,backgrounds,calc,matrix}
\usepackage{pgfplots}
\pgfplotsset{compat=1.17}
\usepgfplotslibrary{groupplots}

% ---------- Source-code listings ----------
\usepackage{listings}

% ---------- Headers / footers ----------
\usepackage{fancyhdr}
\pagestyle{fancy}
\fancyhf{}
\fancyfoot[C]{\thepage}
\renewcommand{\headrulewidth}{0pt}
\fancypagestyle{plain}{%
  \fancyhf{}%
  \fancyfoot[C]{\thepage}%
  \renewcommand{\headrulewidth}{0pt}%
}

% ---------- Chapter / section heading style ----------
\usepackage[explicit]{titlesec}
\titleformat{\chapter}[display]
  {\normalfont\bfseries\centering}
  {\MakeUppercase{\chaptertitlename}\ \thechapter}{12pt}
  {\Large\MakeUppercase{#1}}
\titlespacing*{\chapter}{0pt}{-20pt}{28pt}
\titleformat{\section}{\normalfont\large\bfseries}{\thesection}{1em}{#1}
\titleformat{\subsection}{\normalfont\normalsize\bfseries}{\thesubsection}{1em}{#1}
\titleformat{\subsubsection}{\normalfont\normalsize\itshape}{\thesubsubsection}{1em}{#1}

% ---------- Table of contents depth ----------
\setcounter{tocdepth}{2}
\setcounter{secnumdepth}{3}

% ---------- Captions ----------
\captionsetup{font=small,labelfont=bf,justification=centering}

% ---------- Hyperlinks (load late) ----------
\usepackage[hidelinks,breaklinks=true]{hyperref}

% ============================================================
% IPA SUPPORT
% The phoneme inventory used in this thesis is written directly as
% Unicode in the source; each symbol is mapped onto its `tipa`
% equivalent so the document compiles with plain pdfLaTeX.
% ============================================================
% TIPA's *named* symbol commands are used throughout (rather than its
% single-character shortcut codes), because the names are unambiguous and
% stable across TIPA versions.
\DeclareUnicodeCharacter{0251}{\textscripta}       % open back unrounded
\DeclareUnicodeCharacter{0252}{\textturnscripta}   % open back rounded
\DeclareUnicodeCharacter{0254}{\textopeno}         % open-mid back rounded
\DeclareUnicodeCharacter{0259}{\textschwa}         % schwa
\DeclareUnicodeCharacter{025A}{\textrhookschwa}    % rhotic schwa
\DeclareUnicodeCharacter{025B}{\textepsilon}       % open-mid front unrounded
\DeclareUnicodeCharacter{0261}{\textscriptg}       % script g
\DeclareUnicodeCharacter{026A}{\textsci}           % near-close near-front
\DeclareUnicodeCharacter{0279}{\textturnr}         % alveolar approximant
\DeclareUnicodeCharacter{027E}{\textfishhookr}     % alveolar flap
\DeclareUnicodeCharacter{0283}{\textesh}           % voiceless postalveolar fric.
\DeclareUnicodeCharacter{028A}{\textupsilon}       % near-close near-back
\DeclareUnicodeCharacter{028C}{\textturnv}         % open-mid back unrounded
\DeclareUnicodeCharacter{0292}{\textyogh}          % voiced postalveolar fric.
\DeclareUnicodeCharacter{0294}{\textglotstop}      % glottal stop
\DeclareUnicodeCharacter{014B}{\ng}                % velar nasal
\DeclareUnicodeCharacter{03B8}{\texttheta}         % voiceless dental fricative
\DeclareUnicodeCharacter{00E6}{\ae}                % ash
\DeclareUnicodeCharacter{02C8}{\textprimstress}    % primary stress
\DeclareUnicodeCharacter{02CC}{\textsecstress}     % secondary stress
\DeclareUnicodeCharacter{02D0}{\textlengthmark}    % length mark

% Convenience wrapper so phonemes stand out consistently in running text.
\newcommand{\ph}[1]{\mbox{/#1/}}

% ============================================================
% BREAKABLE MONOSPACE
% \texttt never hyphenates, so long model names, file paths and
% environment-variable names overflow the measure. \pth is a
% url-style command: same typewriter look, but it may break after
% "/", "-", "_" and ".". Its argument is verbatim, so underscores
% are written literally (\pth{CLARIO_ERROR_THRESHOLD}).
% Not usable inside a caption or other moving argument.
% ============================================================
\usepackage{url}
\DeclareUrlCommand\pth{\urlstyle{tt}}
\def\UrlBreaks{\do\/\do\-\do\_\do\.\do\+\do\=\do\:}
\setlength{\Urlmuskip}{0mu plus 1mu}

% ============================================================
% LISTINGS STYLE (code appendices)
% Code excerpts are kept ASCII-only so the listing engine never has to
% typeset a Unicode glyph.
% ============================================================
\definecolor{codegray}{gray}{0.95}
\definecolor{codegreen}{rgb}{0.0,0.45,0.15}
\definecolor{codeblue}{rgb}{0.10,0.20,0.65}
\definecolor{codemauve}{rgb}{0.45,0.10,0.45}

\lstdefinestyle{thesisstyle}{
  backgroundcolor=\color{codegray},
  basicstyle=\ttfamily\scriptsize,
  keywordstyle=\color{codeblue}\bfseries,
  commentstyle=\itshape\color{codegreen},
  stringstyle=\color{codemauve},
  numberstyle=\tiny\color{gray},
  numbers=left,
  numbersep=8pt,
  breaklines=true,
  breakatwhitespace=false,
  captionpos=b,
  frame=single,
  framesep=4pt,
  showstringspaces=false,
  columns=fullflexible,
  keepspaces=true,
  tabsize=4,
  upquote=true,
}
\lstset{style=thesisstyle}
\renewcommand{\lstlistingname}{Listing}

% ============================================================
% SHARED TIKZ STYLES FOR THE DESIGN DIAGRAMS
% ============================================================
\tikzset{
  block/.style   = {rectangle, rounded corners=2pt, draw=black!70, fill=blue!5,
                    text width=2.9cm, align=center, minimum height=1cm,
                    font=\small, inner sep=4pt},
  wideblock/.style = {block, text width=4.2cm},
  data/.style    = {rectangle, draw=black!70, fill=green!6, text width=2.9cm,
                    align=center, minimum height=0.9cm, font=\small},
  decision/.style= {diamond, aspect=2, draw=black!70, fill=orange!12,
                    align=center, font=\scriptsize, inner sep=1pt},
  actor/.style   = {circle, draw=black!70, fill=gray!10, minimum size=0.9cm,
                    font=\scriptsize, align=center},
  arrow/.style   = {-{Stealth[length=2.6mm]}, thick, draw=black!70},
  dashedarrow/.style = {-{Stealth[length=2.6mm]}, thick, dashed, draw=black!60},
  lbl/.style     = {font=\scriptsize, midway, fill=white, inner sep=1.5pt},
}

% ============================================================
% MISC
% ============================================================
\newcommand{\clario}{\textsc{Clario}}
\renewcommand{\arraystretch}{1.15}
\interfootnotelinepenalty=10000
\hyphenpenalty=2000
\tolerance=1500
